\subsection*{A cellular prior teaches inference from labelled context}

% P-R01 | fig1__facts: simulator.*; notes/result_summary.md qualitative anchors
We constructed a procedural prior over single-cell prediction problems, varying biological differences together with the conditions under which they are observed. The generator samples hierarchical cell identities, distributed expression programs, rare markers, correlated variation and donor-dependent states. Donor composition, study and batch shifts, sequencing depth and overdispersed count sampling produce heterogeneous observations of these latent signals. Different labels and support budgets turn each simulated dataset into multiple inference problems. This makes experimental variation part of the training distribution, with known task labels and controllable nuisance factors. Aggregate real-data statistics informed the prior's ranges, whereas all examples used to optimize CellPFN's weights were synthetic.

% P-R02 | fig1__facts: episode/input/architecture; current neural-only contract
Each problem comprises labelled support cells and unlabelled queries. Genes retain their correspondence between support and query but have no persistent identity across problems. CellPFN compares each query gene with support-derived class summaries, pools the resulting relations through permutation-invariant statistics and integrates this evidence into class scores. This use of unordered features follows the principle of set-based learning\cite{zaheer2017deepsets}, while preserving the within-task correspondence required for expression comparisons. A covariance estimate fitted to the support cells supplies an additional statistical channel when support is sufficient. The current model produces direct neural predictions; no analytic classifier supplies its starting logits. Prediction on a new task requires labelled context but no gradient updates, and queries do not condition on other query cells.

% P-R03 | R01; T01-T07 | external numerator: qu2026tabiclv2 section4.1
The reference model has 12.33 million trainable parameters and was trained over 384,000 scheduled episodes drawn from 12,000 simulated dataset banks. Reused banks supplied distinct task views; 276,828 episodes carried nonzero loss weight. For scale, TabICLv2 reports approximately 35 million synthetic pretraining datasets\cite{qu2026tabiclv2}, about 90 times CellPFN's scheduled episode count. This comparison describes training-task exposure across different priors and architectures; it does not measure a matched-compute advantage. The resulting question is whether this limited exposure to cellular tasks suffices for useful inference on real cells.

\subsection*{Compact CellPFN is competitive across biological inference regimes}

% P-R04 | fig4 panel definitions; R02-R12; processed Zebra primary view
We evaluated coarse and fine cell-type annotation in peripheral blood mononuclear cells (PBMCs)\cite{hao2021multimodal}, together with cortex annotation within the development collection and on queries from held-out studies in ZEBRA\cite{flotho2024zebra}. Each panel used one, ten or 100 labelled support cells per class. Comparators included support-trained statistical classifiers, released tabular models, and logistic classifiers fitted to Geneformer and scGPT embeddings\cite{theodoris2023geneformer,cui2024scgpt,qu2026tabiclv2,nvidia2026kumo,jaeger2026tabpfn35}. The main comparisons used the same three support draws. PBMC inputs were counts for CellPFN and normalized expression for tabular comparators; the cortex analysis used the exported processed-expression view. Model-specific input routes were retained for the single-cell foundation models.

% P-R05 | R02-R07: equal-weight descriptive averages across12 task/budget cells
Across the four panels and three support sizes, CellPFN using all available genes reached a mean balanced accuracy of 72.9\%. Released TabICLv2 and Kumo large reached 66.8\% and 70.1\%, respectively, using 2,000 highly variable genes, while cellular continued training brought Kumo large to 72.4\%. Geneformer and scGPT with logistic classifiers reached 74.9\% and 69.4\% through their native input routes. These descriptive averages place the compact cellular model in a competitive range without establishing statistical equivalence or a universal ranking.

% P-R06 | R08-R12; FIGURES.md fig4 input/pretraining caveats; O01
The ranking depended on the biological task and the amount of labelled context. For fine PBMC annotation with one support cell per class, Geneformer reached 68.8\% balanced accuracy versus 51.1\% for CellPFN. For cortex queries from held-out studies with ten support cells per class, CellPFN reached 91.9\%, compared with 87.4\% for Geneformer and 87.2\% for scGPT. The foundation-model comparisons combine different pretraining exposure, preprocessing and classifier settings. Moreover, model variants were selected on these development panels, and each training arm currently has one training seed. \unresolvedclaim{O01}{A frozen evaluation on additional, independently held-out cohorts will determine how consistently this competitive performance generalizes.}

\subsection*{Gene-anonymous inference extends to transcriptome-scale inputs}

% P-R07 | R41-R44; R13,R14,R02; fig2 coherent three-draw comparison
We next asked whether CellPFN could use more genes than it encountered during training. The largest training tasks contained 8,192 genes, whereas the PBMC and processed cortex evaluations contained 33,538 and 21,158 genes, respectively. The same shared per-gene computations and symmetric aggregation accepted these larger inputs without expanding a learned gene vocabulary. Across the matched benchmark draws, balanced accuracy increased from 71.2\% with 2,000 highly variable genes to 73.8\% with 8,192, and was 72.9\% with all genes. Additional genes therefore remained usable, although retaining every gene did not maximize aggregate accuracy and was not beneficial in every task.

% P-R08 | R03,R04,R15,R16,R19; runtime/configuration claim rather than impossibility
Released TabICLv2 declined from 66.8\% to 48.3\% and Kumo large from 70.1\% to 66.4\% when moving from 2,000 highly variable genes to all genes. TabPFN-3.5, evaluated under its non-commercial research terms, refused the all-gene cases at its default feature guard\cite{jaeger2026tabpfn35,priorlabs2026tabpfn35card}. These outcomes describe the tested inference configurations, including their column handling; they do not establish a fundamental limit of tabular architectures. CellPFN's advantage here is the ability to retain transcriptome-scale input under its existing inference procedure.

\subsection*{The cellular prior also improves pretrained tabular models}

% P-R09 | T08-T10; R20-R25 | no claim that adaptation tasks are total lifetime exposure
To test whether the prior's utility extends beyond CellPFN, we continued training released TabICLv2, Kumo small and Kumo large on the cellular task stream. Each model received 40,000 additional cellular tasks and 10,000 generic replay tasks over 2,500 updates. With 2,000 highly variable genes at evaluation, one-shot balanced accuracy improved by 7.3, 10.8 and 5.6 percentage points, respectively, averaged across the four core panels. Changes at 100 support cells per class were smaller than one percentage point in magnitude. Thus the largest benefits arose where labelled biological context was most limited.

% P-R10 | R26-R27; fig3 matched draws, unmatched table dimensions
A generic-table continued-training control for Kumo small changed mean balanced accuracy by approximately zero across the eight low-support panel--budget combinations, compared with a 6.1-point gain for cellular continued training. This comparison supports a contribution from the training distribution, although the controls differed in rows and columns per update and equivalent controls were unavailable for the other architectures. The experiments therefore demonstrate two uses of the cellular prior: training a compact specialist from scratch and adapting already pretrained tabular learners. They do not require that the same architecture be optimal for both uses.

\subsection*{Statistical structure helps selectively rather than uniformly}

% P-R11 | fig6 preregistration and October13 none/diag/full refresh
We examined whether explicitly exposing cross-gene structure accounted for CellPFN's performance. The support-fitted covariance channel supplied both per-gene variance information and low-rank relationships between genes. Later comparisons of separately trained models with no operator, diagonal information alone or the full operator showed gains in several support-rich settings. However, the prespecified operator test did not support its primary claim: an operator-attributable loss exceeded the allowed threshold on a primary task. The later descriptive improvements therefore do not establish the operator as uniformly beneficial, and variation across independent training runs remains unmeasured.

% P-R12 | fig6 taskmix; fig7 realism | qualified constructive boundary
Changing which simulated tasks received training weight also produced a tradeoff. Continuation on tasks selected as learnable improved some low-support accuracies but worsened predictive log loss at larger support sizes; the comparison additionally changed the number of cells per update. Simulator diagnostics likewise showed partial agreement with real data rather than uniform fidelity: prominent marker effects were reproduced more closely than the difficulty of few-shot discrimination. Together, these results motivate judging the prior through transfer and targeted controls, alongside distributional checks, rather than treating greater simulated realism or a stronger statistical channel as sufficient evidence of better inference.

\subsection*{Query-gene deletion explains predictions efficiently}

% P-R13 | fig5 contract and comparison with IG; R28-R31
Gene anonymity still permits gene-level explanations because feature names can be restored after prediction. We developed Deep-Set leave-one-out (DS-LOO), which estimates the change in a query's class margin when an individual gene is set to zero, with the support set and fitted covariance held fixed. In contrast to gradient attribution along an input path\cite{sundararajan2017ig}, the target is this finite query-side deletion. DS-LOO uses the pooled gene statistics and a linearization of the prediction head, with corrections for changes in depth and detected-gene fraction. It is an approximation to a model intervention, not a simulated biological gene knockout.

% P-R14 | R28-R35 | fidelity versus cost; exact-all-gene hours are extrapolated
Against exact deletions in six evaluated cases and two model states, the case-level medians of within-cell rank correlation ranged from 0.979 to 0.998. For the reference model, corrected magnitude errors ranged from 0.5\% to 4.9\%. These checks used selected genes and correctly classified cells in full precision. Under matched load, computing all-gene explanations cost approximately 10--13 single-gene exact forward passes. For the PBMC example, DS-LOO took 67.9 seconds; the corresponding 61-hour exhaustive-deletion runtime was extrapolated from individual forward timings rather than measured directly. Shared-GPU timings remain provisional.

% P-R15 | R36,R37,R46,R47 | fixed original runner-up margin, not multiclass error
The explanations recovered recognizable cell-type markers despite the absence of gene-name embeddings: IL7R ranked second for CD4 T cells and NKG7 seventh for natural killer cells in the evaluated class-margin rankings. Removing each PBMC query's 50 most influential genes reversed its margin against the original runner-up class in 45.0\% of evaluated cells, compared with 0.3\% for random expressed genes. This supports faithfulness to the classifier's decisions. It does not show that the same deletions would change biological identity or predict experimental perturbation responses.

\subsection*{A provisional application examines neuropathology-associated cell states}

% P-R16 | R38-R40; supp_rosmap status; Bennett is cohort context only
\unresolvedclaim{O02}{We extended the framework to donor-level neuropathology contrasts in ROSMAP-derived single-nucleus data\cite{bennett2018rosmap}. In the exported analysis with case and control support matched within sex, microglial predictions distinguished the CERAD 4-versus-1 donor groups with an AUC of 0.745 and a reported corrected repeated-split interval of 0.700--0.790. This result remains provisional pending completion of the donor-label audit and a frozen rerun. It measures association with neuropathology, not prospective diagnosis or clinical utility.}

% P-R17 | no completed disease gene panel or independent replication exists
\unresolvedclaim{O03}{Applying query-gene deletion to the corrected disease analysis and testing the findings in independent cohorts will establish whether this approach can connect donor-level prediction to reproducible cellular expression signatures.}
